% ৪.১১ ওয়েব পেজ ডিজাইন পদ্ধতি
\section{ওয়েব পেজ ডিজাইন পদ্ধতি}

\begin{learningbox}
এই টপিক শেষে শিক্ষার্থী পারবে:
\begin{itemize}
\item web page design workflow লিখতে।
\item wireframe, prototype ও mockup আলাদা করতে।
\item একটি simple web page-এর HTML structure পরিকল্পনা করতে।
\end{itemize}
\end{learningbox}

\subsection{Design workflow}
একটি page সরাসরি code করার আগে কী content থাকবে, কোন অংশ কোথায় থাকবে এবং ব্যবহারকারী কীভাবে page ব্যবহার করবে, তা পরিকল্পনা করতে হয়। এতে ভুল কমে এবং page maintain করা সহজ হয়।

\begin{keypointbox}{ধাপ}
\begin{itemize}
\item Requirement: page-এর উদ্দেশ্য ও target user নির্ধারণ।
\item Sitemap: page সম্পর্ক ও navigation নির্ধারণ।
\item Wireframe: layout-এর rough sketch।
\item Prototype/mockup: visual design পরীক্ষা।
\item HTML/CSS coding: structure ও style তৈরি।
\item Testing: browser, link, form, mobile layout যাচাই।
\end{itemize}
\end{keypointbox}

\begin{center}
\fbox{\parbox{0.9\textwidth}{
\centering
\textbf{Wireframe placeholder: Student profile page}\\[4pt]
Header with title\\
Profile image + short bio\\
Education table\\
Contact form\\
Footer
}}
\end{center}

\begin{examplebox}{Mini page plan}
``আমার কলেজ'' page-এ heading, কলেজের সংক্ষিপ্ত পরিচিতি, একটি campus image, notice link, একটি routine table এবং contact form রাখা যেতে পারে। এতে Chapter 4-এর প্রায় সব core tag একসাথে practice হয়।
\end{examplebox}

\begin{boardbox}{বোর্ড ফোকাস}
প্রশ্নে web page design চাইলে শুধু code নয়, page-এর expected output বা layout ধারণাও লিখলে উত্তর ভালো হয়।
\end{boardbox}

\begin{practicebox}
\begin{enumerate}
\item Wireframe কী?
\item Prototype ও final website-এর পার্থক্য লেখ।
\item একটি school homepage-এর পাঁচটি অংশের নাম লেখ।
\end{enumerate}
\end{practicebox}
